\chapter{\eh{handshake} Capa 5 --- Sesión (Session Layer)}

\begin{concepto}
La capa de sesión establece, gestiona y termina \textbf{sesiones} de
comunicación entre aplicaciones. Se encarga de sincronizar el diálogo,
controlar quién habla cuándo, y manejar puntos de control (checkpoints)
para reanudar sesiones interrumpidas.

\emoji{warning} \textbf{Nota sobre clasificación:} SSH, NetBIOS y otros protocolos
que ``viven'' en capa 5 son a veces clasificados en capa 7 por otras fuentes.
En la práctica real, los protocolos no siguen OSI perfectamente. Lo importante
es entender qué hace cada protocolo, no en qué caja exacta meterlo.
\end{concepto}

\section{\eh{locked-with-key} SSH --- Secure Shell (RFC 4251--4254)}

\begin{concepto}
\textbf{SSH} proporciona acceso remoto seguro a sistemas. Todo el tráfico
está \textbf{cifrado}: contraseñas, comandos, y respuestas. Reemplazó
completamente a Telnet y rlogin.

\emoji{light-bulb} \textbf{Analogía:} Telnet es como hablar en voz alta
en un café lleno de gente --- cualquiera escucha tu contraseña.
SSH es como estar en una cabina de sonido insonorizada y con candado:
solo tú y el servidor saben qué se dijo.
\end{concepto}

\subsection{SSH vs Telnet}

\begin{table}[h!]
\centering
\renewcommand{\arraystretch}{1.4}
\begin{tabular}{lll}
\toprule
\textbf{Característica} & \textbf{Telnet} & \textbf{SSH} \\
\midrule
Puerto          & 23              & 22 \\
Cifrado         & \emoji{cross-mark} Ninguno & \emoji{check-mark-button} Sí (AES, ChaCha20) \\
Autenticación   & Contraseña en texto claro & Clave pública o contraseña cifrada \\
Integridad      & \emoji{cross-mark} No    & \emoji{check-mark-button} HMAC \\
Seguridad       & \emoji{skull} Inseguro   & \emoji{shield} Seguro \\
Estado actual   & Obsoleto, no usar         & Estándar de la industria \\
\bottomrule
\end{tabular}
\caption{Comparativa SSH vs Telnet}
\end{table}

\subsection{Arquitectura SSH (3 protocolos en uno)}

SSH se compone de 3 sub-protocolos apilados:
\begin{enumerate}
    \item \textbf{SSH Transport Layer (RFC 4253):} establece la conexión TCP,
          negocia algoritmos y realiza el key exchange (Diffie-Hellman).
    \item \textbf{SSH Authentication Protocol (RFC 4252):} autentica al usuario
          con contraseña o clave pública.
    \item \textbf{SSH Connection Protocol (RFC 4254):} multiplexa múltiples
          canales sobre la conexión segura (sesión de terminal, port forwarding, SFTP).
\end{enumerate}

\subsection{Proceso de conexión SSH}

\begin{lstlisting}
  Cliente                         Servidor
     │──── TCP SYN ──────────────►│
     │◄─── TCP SYN-ACK ───────────│
     │──── TCP ACK ──────────────►│
     │                            │
     │◄─── SSH-2.0-OpenSSH_x.x ──│  (identificación de versión)
     │──── SSH-2.0-OpenSSH_x.x ──►│
     │                            │
     │◄─── Server Key Exchange ───│  (algoritmos + Diffie-Hellman params)
     │──── Client Key Exchange ──►│  (valor DH público del cliente)
     │                            │  (ambos calculan el shared secret)
     │◄─── New Keys ──────────────│  (todo lo siguiente va cifrado)
     │──── New Keys ─────────────►│
     │                            │
     │──── [ENC] auth request ───►│  (usuario + método de auth)
     │◄─── [ENC] auth success ────│
     │                            │
     │──── [ENC] channel open ───►│  (abrir sesión de terminal)
     │◄─── [ENC] channel conf ────│
     │═══════ Sesión cifrada ══════│
\end{lstlisting}

\subsection{Autenticación por clave pública}

\begin{concepto}
Más segura que contraseña porque la clave privada \textit{nunca} sale de
tu máquina. Basada en criptografía asimétrica: lo que una clave cifra,
solo la otra puede descifrar.
\end{concepto}

\begin{enumerate}
    \item Generas un par de claves: \texttt{ssh-keygen -t ed25519 -C "tu@email.com"}
    \item Clave pública (\texttt{\textasciitilde/.ssh/id\_ed25519.pub}) → va al servidor
          en \texttt{\textasciitilde/.ssh/authorized\_keys}
    \item Al conectar: el servidor cifra un desafío aleatorio con tu clave pública
    \item Tu cliente lo descifra con la clave privada y responde
    \item Sin contraseña que viajar por la red → sin riesgo de captura
\end{enumerate}

\begin{table}[h!]
\centering
\renewcommand{\arraystretch}{1.3}
\begin{tabular}{lll}
\toprule
\textbf{Algoritmo} & \textbf{Bits} & \textbf{Nota} \\
\midrule
RSA       & 2048--4096  & Antiguo, funciona bien pero lento \\
ECDSA     & 256--521    & Más rápido que RSA, elliptic curve \\
Ed25519   & 256 (equiv) & Recomendado hoy: más rápido, más seguro \\
\bottomrule
\end{tabular}
\caption{Algoritmos de clave pública en SSH}
\end{table}

\subsection{Port Forwarding con SSH}

\begin{codigo}
\begin{lstlisting}
# Local forwarding: acceder puerto remoto como si fuera local
# Acceder MySQL (puerto 3306) de servidor remoto como localhost:3307
ssh -L 3307:localhost:3306 usuario@servidor.com

# Remote forwarding: exponer puerto local hacia el servidor remoto
ssh -R 8080:localhost:3000 usuario@servidor.com

# Dynamic forwarding: SSH como proxy SOCKS5
ssh -D 1080 usuario@servidor.com
\end{lstlisting}
\end{codigo}

\section{\eh{left-right-arrow} NetBIOS y SMB}

\begin{concepto}
\textbf{NetBIOS} (Network Basic Input/Output System) proporciona servicios
de nombres y sesiones en redes Windows locales. No es routable por IP
--- solo funciona en la LAN.

\textbf{SMB} (Server Message Block) usa NetBIOS (o directamente TCP/IP en SMB2+)
para compartir archivos, impresoras y recursos en redes Windows.
Puerto: 445 (SMB directo) o 139 (SMB sobre NetBIOS).
\end{concepto}
